%=============================================================
%  PEDAGOGICAL SUPPORT MATERIAL
%  Quantum Mechanics & Quantum Chemistry Programme
%  Learning Tiers + Assessment Artifact Set (Revised Taxonomy)
%=============================================================
\documentclass[11pt,a4paper]{article}

%--- Packages ------------------------------------------------
\usepackage[T1]{fontenc}
\usepackage[utf8]{inputenc}
\usepackage{fix-cm}
\usepackage[margin=2.5cm]{geometry}
\usepackage{amsmath,amssymb}
\usepackage{booktabs}
\usepackage{array}
\usepackage{longtable}
\usepackage{tabularx}
\usepackage{multirow}
\usepackage{enumitem}
\usepackage{sectsty}
\usepackage{fancyhdr}
\usepackage[dvipsnames,svgnames,table]{xcolor}
\usepackage{hyperref}
\usepackage{parskip}
\usepackage{tcolorbox}
\tcbuselibrary{skins,breakable}

%--- Colours -------------------------------------------------
\definecolor{darkblue}{RGB}{10,40,90}
\definecolor{midblue}{RGB}{30,80,160}
\definecolor{lightblue}{RGB}{220,233,255}
\definecolor{accentgold}{RGB}{180,140,20}
\definecolor{lightgold}{RGB}{255,248,220}
\definecolor{darkgray}{RGB}{60,60,60}
\definecolor{lightgray}{RGB}{245,245,245}
% Tier colours
\definecolor{tier1col}{RGB}{180,230,180}
\definecolor{tier2col}{RGB}{150,210,240}
\definecolor{tier3col}{RGB}{200,180,240}
\definecolor{tier4col}{RGB}{255,210,150}
\definecolor{tier5col}{RGB}{255,170,150}
% Undergrad track colours
\definecolor{ugblue}{RGB}{25,90,170}
\definecolor{lightugblue}{RGB}{215,230,255}
\definecolor{uggreen}{RGB}{20,120,70}
\definecolor{lightuggreen}{RGB}{210,245,225}
% Grad track colours
\definecolor{gradred}{RGB}{160,30,30}
\definecolor{lightgradred}{RGB}{255,220,215}
\definecolor{gradorange}{RGB}{170,80,10}
\definecolor{lightgradorange}{RGB}{255,235,210}
% Research track colours
\definecolor{respath}{RGB}{90,20,130}
\definecolor{lightrespath}{RGB}{240,220,255}
% Rubric colour
\definecolor{artpurple}{RGB}{100,30,130}
\definecolor{lightartpurple}{RGB}{235,215,245}

%--- Section fonts -------------------------------------------
\sectionfont{\Large\bfseries\color{darkblue}\sectionrule{0ex}{0pt}{-1ex}{1.5pt}}
\subsectionfont{\large\bfseries\color{midblue}}
\subsubsectionfont{\normalsize\bfseries\color{darkgray}}

%--- Hyperref ------------------------------------------------
\hypersetup{
  colorlinks=true,
  linkcolor=midblue,
  urlcolor=midblue,
  citecolor=midblue,
  pdftitle={Pedagogical Support Material -- QM Programme},
}

%--- Header / Footer -----------------------------------------
\pagestyle{fancy}
\fancyhf{}
\fancyhead[L]{\small\color{darkgray}\textit{QM Programme --- Pedagogical Support}}
\fancyhead[R]{\small\color{darkgray}\thepage}
\fancyfoot[C]{\small\color{darkgray}\textit{For instructors and students}}
\renewcommand{\headrulewidth}{0.4pt}

%--- tcolorbox styles ----------------------------------------
\tcbset{
  tierbox/.style={
    enhanced, breakable,
    colframe=#1!80!black,
    colback=#1!15!white,
    fonttitle=\bfseries\color{white}, coltitle=white,
    attach boxed title to top left={yshift=-2mm, xshift=4mm},
    boxed title style={colback=#1!80!black, sharp corners, boxrule=0.5pt},
    sharp corners=south, boxrule=0.8pt,
    left=6pt, right=6pt, top=8pt, bottom=6pt,
  },
  % Undergrad concept/problem boxes
  ugconceptbox/.style={
    enhanced, breakable,
    colback=lightugblue, colframe=ugblue,
    fonttitle=\bfseries\color{white}, coltitle=white,
    attach boxed title to top left={yshift=-2mm, xshift=4mm},
    boxed title style={colback=ugblue, colframe=ugblue!80!black,
      sharp corners, boxrule=0.5pt},
    sharp corners=south, boxrule=0.8pt,
    left=6pt, right=6pt, top=8pt, bottom=6pt,
  },
  ugprobbox/.style={
    enhanced, breakable,
    colback=lightuggreen, colframe=uggreen,
    fonttitle=\bfseries\color{white}, coltitle=white,
    attach boxed title to top left={yshift=-2mm, xshift=4mm},
    boxed title style={colback=uggreen, colframe=uggreen!80!black,
      sharp corners, boxrule=0.5pt},
    sharp corners=south, boxrule=0.8pt,
    left=6pt, right=6pt, top=8pt, bottom=6pt,
  },
  % Grad concept/problem boxes
  gradconceptbox/.style={
    enhanced, breakable,
    colback=lightgradred, colframe=gradred,
    fonttitle=\bfseries\color{white}, coltitle=white,
    attach boxed title to top left={yshift=-2mm, xshift=4mm},
    boxed title style={colback=gradred, colframe=gradred!80!black,
      sharp corners, boxrule=0.5pt},
    sharp corners=south, boxrule=0.8pt,
    left=6pt, right=6pt, top=8pt, bottom=6pt,
  },
  gradprobbox/.style={
    enhanced, breakable,
    colback=lightgradorange, colframe=gradorange,
    fonttitle=\bfseries\color{white}, coltitle=white,
    attach boxed title to top left={yshift=-2mm, xshift=4mm},
    boxed title style={colback=gradorange, colframe=gradorange!80!black,
      sharp corners, boxrule=0.5pt},
    sharp corners=south, boxrule=0.8pt,
    left=6pt, right=6pt, top=8pt, bottom=6pt,
  },
  % Research track
  researchbox/.style={
    enhanced, breakable,
    colback=lightrespath, colframe=respath,
    fonttitle=\bfseries\color{white}, coltitle=white,
    attach boxed title to top left={yshift=-2mm, xshift=4mm},
    boxed title style={colback=respath, colframe=respath!80!black,
      sharp corners, boxrule=0.5pt},
    sharp corners=south, boxrule=0.8pt,
    left=6pt, right=6pt, top=8pt, bottom=6pt,
  },
  % Rubric and utility boxes
  rubricbox/.style={
    enhanced,
    colback=lightartpurple, colframe=artpurple,
    fonttitle=\bfseries\color{white}, coltitle=white,
    attach boxed title to top left={yshift=-2mm, xshift=4mm},
    boxed title style={colback=artpurple, sharp corners, boxrule=0.5pt},
    sharp corners=south, boxrule=0.7pt,
    left=5pt, right=5pt, top=6pt, bottom=5pt,
  },
  goldbox/.style={
    enhanced, breakable,
    colback=lightgold, colframe=accentgold,
    boxrule=0.8pt,
    left=6pt, right=6pt, top=6pt, bottom=6pt,
    fontupper=\small,
  },
  notebox/.style={
    enhanced,
    colback=lightblue, colframe=midblue,
    boxrule=0.6pt, rounded corners,
    left=5pt, right=5pt, top=4pt, bottom=4pt,
    fontupper=\small\itshape,
  },
  trackheaderUG/.style={
    enhanced,
    colback=ugblue!12!white, colframe=ugblue,
    boxrule=1.2pt, sharp corners,
    left=6pt, right=6pt, top=5pt, bottom=5pt,
    fontupper=\bfseries\large\color{ugblue},
  },
  trackheaderGRAD/.style={
    enhanced,
    colback=gradred!12!white, colframe=gradred,
    boxrule=1.2pt, sharp corners,
    left=6pt, right=6pt, top=5pt, bottom=5pt,
    fontupper=\bfseries\large\color{gradred},
  },
  trackheaderRES/.style={
    enhanced,
    colback=respath!12!white, colframe=respath,
    boxrule=1.2pt, sharp corners,
    left=6pt, right=6pt, top=5pt, bottom=5pt,
    fontupper=\bfseries\large\color{respath},
  },
}

\newcommand{\instructornote}[1]{%
  \begin{tcolorbox}[notebox]\textbf{Instructor note:} #1\end{tcolorbox}%
}
\newcommand{\tiertag}[2]{%
  \colorbox{#1}{\small\bfseries\color{darkblue}\ #2\ }%
}
\newcommand{\essential}{$\bigstar\bigstar\bigstar$}
\newcommand{\recommended}{$\bigstar\bigstar$}
\newcommand{\usefulsup}{$\bigstar$}

%=============================================================
\begin{document}
%=============================================================

%--- Title Page ----------------------------------------------
\begin{titlepage}
\pagecolor{darkblue}
\color{white}
\vspace*{2.5cm}
\begin{center}
  {\fontsize{26}{32}\selectfont\bfseries Pedagogical Support Material}\\[0.5cm]
  {\fontsize{16}{22}\selectfont Quantum Mechanics \& Quantum Chemistry Programme}\\[0.6cm]
  {\color{accentgold}\rule{10cm}{2pt}}\\[0.6cm]
  {\large\itshape For Students and Instructors}\\[2cm]
  \begin{tabular}{cl}
    \textbf{Section~1:} & Learning Tiers (HS $\to$ PhD)\\[2pt]
                        & \quad Definitions, Competencies \& Worked Example (Module~I.2)\\[8pt]
    \textbf{Section~2:} & Assessment Artifact Set --- Revised Taxonomy\\[2pt]
                        & \quad \textit{Undergraduate Track} (SCQU, CCQU, SPSU, CPSU, SPJU, CPJU)\\[2pt]
                        & \quad \textit{Graduate Track} (SCQG, CCQG, SPSG, CPSG, SPJG, CPJG)\\[2pt]
                        & \quad \textit{Research Track} (WRQ $\times$5, ORQ $\times$5)\\[8pt]
    \textbf{Section~3:} & Tiered Bibliography --- Structure \& Design Principles\\[2pt]
                        & \quad (Full reading lists in companion document\\[2pt]
                        & \quad \textit{QM Programme: Tiered Bibliography})\\
  \end{tabular}\\[2cm]
  {\color{accentgold}\rule{10cm}{1pt}}\\[0.5cm]
  {\small\itshape v3 --- Companion to the Course Syllabus.\\
  Sets the shared language for \emph{how} material is taught and assessed\\
  across all Series and Modules of the programme.\\[4pt]
  See also: \textbf{QM Programme: Tiered Bibliography} (separate companion document)}
\end{center}
\vfill
\end{titlepage}
\nopagecolor

\tableofcontents
\newpage

%=============================================================
\section*{Introduction}
\addcontentsline{toc}{section}{Introduction}
%=============================================================

\begin{tcolorbox}[goldbox]
This document serves two audiences simultaneously.

\medskip\noindent\textbf{For students:} It explains what is expected at each stage of your
education and what kinds of tasks you will be asked to perform. Understanding the tier
system and the artifact taxonomy before a course begins removes ambiguity about difficulty
levels and grading expectations.

\medskip\noindent\textbf{For instructors:} It provides a reusable, course-independent
specification for how to differentiate content and assessment. The tier system and
revised artifact set described here apply uniformly across \emph{every lecture} of
every module in the programme.
\end{tcolorbox}

\medskip
\noindent\textbf{How the two sections connect.}
Section~1 defines the \emph{audience} dimension. Section~2 defines the
\emph{assessment} dimension, now organised into three tracks:

\begin{itemize}[leftmargin=1.8em]
  \item \textcolor{ugblue}{\textbf{Undergraduate track}} (Tiers~1--3): six artifact types,
        each with a simple and complex variant labelled \texttt{U}
  \item \textcolor{gradred}{\textbf{Graduate track}} (Tiers~4--5): six parallel artifact
        types with deeper scope, labelled \texttt{G}
  \item \textcolor{respath}{\textbf{Research track}} (Tiers~4--5): five well-defined
        and five open-ended research questions
\end{itemize}

\newpage
%#############################################################
\section{Learning Tiers}
\label{sec:tiers}
%#############################################################

\begin{tcolorbox}[goldbox]
\textbf{Overview.}
The programme is designed for a \emph{mixed cohort}: a lecture may be attended
simultaneously by students with only secondary-school mathematics and PhD
researchers extending their foundations.
The five-tier system ensures every student finds appropriate entry points and
stretch goals without fragmenting the cohort into separate tracks.

\medskip\noindent The tiers are \emph{not grade labels}. They describe the prerequisite
competency set a student brings. A student may operate at different tiers for
different topics (e.g.\ Tier~3 for linear algebra, Tier~2 for group theory).

\medskip\noindent\textbf{Track alignment:}
Tiers~1--3 correspond to the \textcolor{ugblue}{\textbf{Undergraduate track}};
Tiers~4--5 correspond to the \textcolor{gradred}{\textbf{Graduate track}}
and \textcolor{respath}{\textbf{Research track}}.
\end{tcolorbox}

%-------------------------------------------------------------
\subsection{Tier Ladder at a Glance}

\begin{center}
\begin{tabular}{cllp{6cm}}
\toprule
\textbf{Tier} & \textbf{Label} & \textbf{Track} & \textbf{One-line description}\\
\midrule
1 & High School (HS)       & \textcolor{ugblue}{Undergrad}  &
    Amplitudes are complex arrows; probabilities are squared lengths\\
2 & Beginning UG (BegUG)   & \textcolor{ugblue}{Undergrad}  &
    Bases and components; compute probabilities with $\langle n|\psi\rangle$\\
3 & Advanced UG (AdvUG)    & \textcolor{ugblue}{Undergrad}  &
    Inequalities; geometry and internal consistency of the formalism\\
4 & Master's (MSc)         & \textcolor{gradred}{Graduate}  &
    Rays and projective space; symmetries via Wigner's theorem\\
5 & PhD                    & \textcolor{gradred}{Graduate}  &
    Rigorous Hilbert space; separability; domain subtleties\\
\bottomrule
\end{tabular}
\end{center}

%-------------------------------------------------------------
\subsection{Tier Definitions and Competencies}

\begin{tcolorbox}[tierbox=tier1col, title={Tier 1 --- High School (HS)\ \ \normalfont\small[\textcolor{ugblue}{Undergraduate track}]}]
\textbf{Goal:} Build the \emph{intuition tools} needed for amplitudes and interference.

\medskip\noindent\textbf{Core competencies:}
\begin{itemize}[leftmargin=1.5em,itemsep=2pt]
  \item Complex numbers as ``2D arrows'': $a+ib$, magnitude $|z|$, phase $\arg z$
  \item Magnitude squared as probability: $P\propto|A|^2$
  \item Vector addition analogy: adding complex amplitudes before squaring gives interference
  \item Basic probability: events, normalisation (``total probability = 1'')
\end{itemize}

\medskip\noindent\textbf{You should be able to:}
\begin{itemize}[leftmargin=1.5em,itemsep=2pt]
  \item Compute $|a+ib|^2$
  \item Understand why $|A_1+A_2|^2\neq|A_1|^2+|A_2|^2$ in general
\end{itemize}
\end{tcolorbox}

\medskip
\begin{tcolorbox}[tierbox=tier2col, title={Tier 2 --- Beginning Undergraduate (BegUG)\ \ \normalfont\small[\textcolor{ugblue}{Undergraduate track}]}]
\textbf{Goal:} Become fluent in the \emph{working language} of bras, kets, bases,
and probabilities.

\medskip\noindent\textbf{Core competencies:}
\begin{itemize}[leftmargin=1.5em,itemsep=2pt]
  \item Normalisation: $\langle\psi|\psi\rangle=1$
  \item Orthonormal bases: $\langle m|n\rangle=\delta_{mn}$
  \item Components as inner products: $c_n=\langle n|\psi\rangle$
  \item Probabilities: $P(n)=|c_n|^2$
  \item Insert identity (discrete): $\displaystyle\sum_n|n\rangle\langle n|=\hat{\mathbb{1}}$
\end{itemize}

\medskip\noindent\textbf{You should be able to:}
\begin{itemize}[leftmargin=1.5em,itemsep=2pt]
  \item Expand a state in a basis and compute probabilities
  \item Translate between $|\psi\rangle$ and its coordinate column vector in a chosen basis
\end{itemize}
\end{tcolorbox}

\medskip
\begin{tcolorbox}[tierbox=tier3col, title={Tier 3 --- Advanced Undergraduate (AdvUG)\ \ \normalfont\small[\textcolor{ugblue}{Undergraduate track}]}]
\textbf{Goal:} Understand why the formalism is \emph{internally consistent} and geometric.

\medskip\noindent\textbf{Core competencies:}
\begin{itemize}[leftmargin=1.5em,itemsep=2pt]
  \item Cauchy--Schwarz inequality:
        $|\langle\phi|\psi\rangle|^2\le\langle\phi|\phi\rangle\langle\psi|\psi\rangle$
  \item Triangle inequality: $\|\phi+\psi\|\le\|\phi\|+\|\psi\|$
  \item Metric/geometric viewpoint: overlaps as cosines, distance from $\|\phi-\psi\|$
\end{itemize}

\medskip\noindent\textbf{You should be able to:}
\begin{itemize}[leftmargin=1.5em,itemsep=2pt]
  \item Prove basic inequalities in Hilbert space
  \item Explain why overlaps and probabilities are bounded by 1 for normalised states
  \item See superposition as linear structure, not mystery
\end{itemize}
\end{tcolorbox}

\medskip
\begin{tcolorbox}[tierbox=tier4col, title={Tier 4 --- Master's (MSc)\ \ \normalfont\small[\textcolor{gradred}{Graduate track}]}]
\textbf{Goal:} Shift from ``vectors as states'' to \emph{rays as states} and connect
to symmetry theory.

\medskip\noindent\textbf{Core competencies:}
\begin{itemize}[leftmargin=1.5em,itemsep=2pt]
  \item Rays / projective Hilbert space: physical states are equivalence classes
        $|\psi\rangle\sim e^{i\theta}|\psi\rangle$; pure states are points in
        $\mathbb{P}(\mathcal{H})$
  \item Wigner theorem preview: transformations preserving
        $|\langle\phi|\psi\rangle|^2$ are unitary or antiunitary
\end{itemize}

\medskip\noindent\textbf{You should be able to:}
\begin{itemize}[leftmargin=1.5em,itemsep=2pt]
  \item Explain why global phase is unphysical and why rays are the correct objects
  \item State Wigner's theorem and why it matters for symmetries
\end{itemize}
\end{tcolorbox}

\medskip
\begin{tcolorbox}[tierbox=tier5col, title={Tier 5 --- PhD\ \ \normalfont\small[\textcolor{gradred}{Graduate track}]}]
\textbf{Goal:} Make all foundational statements \emph{mathematically precise} and know
where physics shorthand hides subtleties.

\medskip\noindent\textbf{Core competencies:}
\begin{itemize}[leftmargin=1.5em,itemsep=2pt]
  \item Precise definition of Hilbert space (inner product + completeness)
  \item Separability (countable orthonormal basis / countable dense subset)
  \item Rays as 1D subspaces; $\mathbb{P}(\mathcal{H})$ as a metric space
  \item Operator subtleties: Hermitian vs.\ self-adjoint; generalised eigenvectors;
        rigged Hilbert spaces
\end{itemize}

\medskip\noindent\textbf{You should be able to:}
\begin{itemize}[leftmargin=1.5em,itemsep=2pt]
  \item Distinguish informal bra-ket from rigorous operator statements
  \item Discuss why separability matters for spectral decompositions
  \item Frame state space as a projective object and symmetry as structure on that space
\end{itemize}
\end{tcolorbox}

%-------------------------------------------------------------
\subsection{What Changes Across Tiers (Same Artifact, Different Depth)}

\instructornote{The same assignment can be given to all tiers simultaneously.
What changes is the \emph{language}, the required \emph{rigor}, and the \emph{scope}.
Below: ``phase invariance'' across all five tiers.}

\begin{center}
\begin{tabular}{lp{10.5cm}}
\toprule
\textbf{Tier} & \textbf{How ``phase invariance'' is engaged}\\
\midrule
HS     & ``Why doesn't the arrow's overall rotation matter for length?'' --- intuitive analogy\\
BegUG  & Compute $|e^{i\theta}c_n|^2=|c_n|^2$; verify for a 2-component example\\
AdvUG  & Prove phase invariance for all probabilities in any orthonormal basis expansion\\
MSc    & Well-definedness of Born-rule map on rays in $\mathbb{P}(\mathcal{H})$\\
PhD    & Formalise as $U(1)$ fiber structure; Hopf fibration
        $S^1\to S^{2n+1}\to\mathbb{CP}^n$\\
\bottomrule
\end{tabular}
\end{center}

%-------------------------------------------------------------
\subsection{Worked Example: Module I.2 --- Lecture L01}

The following shows how \textbf{Lecture L01: \textit{Why Dirac Notation? States as
Rays \& Superposition}} is differentiated across all five tiers. This pattern repeats
for every lecture in the programme.

\begin{tcolorbox}[goldbox]
\textbf{L01 topic:} States, kets, rays, global phase, Born rule preview.
Dual Sakurai (operational) and Dirac (abstract) perspectives.
\end{tcolorbox}

\begin{tcolorbox}[tierbox=tier1col, title={Tier 1 (HS) --- L01}]
\textbf{Focus:} $|z|^2$ as probability; interference with two paths; ket as a
``direction arrow''; why global rotation of the arrow changes nothing.\\[4pt]
\textbf{Signature problem:} Given $A_1=\tfrac{1}{\sqrt{2}}$, $A_2=\tfrac{i}{\sqrt{2}}$,
compute total probability; verify it equals 1.
\end{tcolorbox}

\medskip
\begin{tcolorbox}[tierbox=tier2col, title={Tier 2 (BegUG) --- L01}]
\textbf{Focus:} Normalisation; components $c_n=\langle n|\psi\rangle$; notation dictionary.\\[4pt]
\textbf{Signature problem:} $|\psi\rangle=\tfrac{3}{5}|{+}\rangle+\tfrac{4}{5}|{-}\rangle$.
Verify normalisation; find $P(+)$ and $P(-)$; multiply by $e^{i\pi/3}$; confirm
probabilities unchanged.
\end{tcolorbox}

\medskip
\begin{tcolorbox}[tierbox=tier3col, title={Tier 3 (AdvUG) --- L01}]
\textbf{Focus:} Prove global phase is exactly unobservable; distinguish global from
relative phase.\\[4pt]
\textbf{Signature problem:} Prove that for any normalised $|\psi\rangle$, any ONB
$\{|n\rangle\}$, and any $\theta\in\mathbb{R}$:
$|\langle n|e^{i\theta}|\psi\rangle|^2=|\langle n|\psi\rangle|^2$ for all $n$.
\end{tcolorbox}

\medskip
\begin{tcolorbox}[tierbox=tier4col, title={Tier 4 (MSc) --- L01}]
\textbf{Focus:} Rays as equivalence classes; projective Hilbert space
$\mathbb{P}(\mathcal{H})$; Born-rule map well-defined on rays; preview of Wigner's theorem.\\[4pt]
\textbf{Signature problem:} Define the ray equivalence relation; show Born-rule probability
depends only on $[\psi]$, not the representative; frame what a ``symmetry'' of physical
states means.
\end{tcolorbox}

\medskip
\begin{tcolorbox}[tierbox=tier5col, title={Tier 5 (PhD) --- L01}]
\textbf{Focus:} Hilbert space completeness; separability; rays as 1D subspaces;
Fubini--Study metric; full statement of Wigner's theorem.\\[4pt]
\textbf{Signature problem:} (a) State the definition of a separable Hilbert space.
(b) Why does separability guarantee a countable ONB?
(c) State Wigner's theorem precisely, including the $U(1)$ lifting ambiguity.
(d) Give one example of an antiunitary symmetry and explain its physical significance.
\end{tcolorbox}

\newpage
%#############################################################
\section{Assessment Artifact Set --- Revised Taxonomy}
\label{sec:artifacts}
%#############################################################

\begin{tcolorbox}[goldbox]
\textbf{Design goals.} Each lecture assesses four capabilities:
\begin{enumerate}[leftmargin=2em,itemsep=2pt]
  \item \textbf{Conceptual understanding} --- can you explain it?
  \item \textbf{Procedural fluency} --- can you compute it?
  \item \textbf{Structural mastery} --- can you prove or derive it?
  \item \textbf{Transfer \& creation} --- can you apply it to a new context or ask new questions?
\end{enumerate}

\medskip
The revised taxonomy provides \textbf{parallel sets} for the undergraduate and
graduate tracks, plus a shared research track, for a total of \textbf{14 artifact
types} per lecture. Each type has a defined count, scope, and rubric.
\end{tcolorbox}

%-------------------------------------------------------------
\subsection{Complete Artifact Taxonomy at a Glance}

\begin{center}
\begin{tabular}{llllr}
\toprule
\textbf{Code} & \textbf{Full Name} & \textbf{Track} & \textbf{Tiers} & \textbf{Count}\\
\midrule
\multicolumn{5}{l}{\textit{\textcolor{ugblue}{\textbf{Undergraduate Concept \& Problem Artifacts}}}}\\
SCQU & Simple Concept Questions --- Undergrad   & UG & 1--3 & 10 questions\\
CCQU & Complex Concept Questions --- Undergrad  & UG & 2--3 & 10 questions\\
SPSU & Simple Problem Set --- Undergrad         & UG & 1--3 & 10 problems\\
CPSU & Complex Problem Set --- Undergrad        & UG & 3    & 5 problems\\
SPJU & Simple Project --- Undergrad             & UG & 2--3 & 1 project\\
CPJU & Complex Project --- Undergrad            & UG & 3    & 1 project\\
\midrule
\multicolumn{5}{l}{\textit{\textcolor{gradred}{\textbf{Graduate Concept \& Problem Artifacts}}}}\\
SCQG & Simple Concept Questions --- Grad        & Grad & 4--5 & 10 questions\\
CCQG & Complex Concept Questions --- Grad       & Grad & 4--5 & 10 questions\\
SPSG & Simple Problem Set --- Grad              & Grad & 4--5 & 10 problems\\
CPSG & Complex Problem Set --- Grad             & Grad & 4--5 & 5 problems\\
SPJG & Simple Project --- Grad                  & Grad & 4--5 & 5 projects\\
CPJG & Complex Project --- Grad                 & Grad & 4--5 & 5 projects\\
\midrule
\multicolumn{5}{l}{\textit{\textcolor{respath}{\textbf{Research Track (shared Tiers 4--5)}}}}\\
WRQ  & Well-Defined Research Question           & Research & 4--5 & 5 questions\\
ORQ  & Open-Ended Research Question             & Research & 4--5 & 5 questions\\
\midrule
 & \multicolumn{3}{l}{\textbf{Total artifacts per lecture}} & \textbf{92 items}\\
\bottomrule
\end{tabular}
\end{center}

%-------------------------------------------------------------
\subsection{Capability--Artifact Mapping}

\begin{center}
\begin{tabularx}{\textwidth}{lXX}
\toprule
\textbf{Capability} & \textbf{UG artifacts} & \textbf{Grad artifacts}\\
\midrule
Conceptual understanding & SCQU, CCQU & SCQG, CCQG\\
Procedural fluency       & SPSU       & SPSG\\
Structural mastery       & CPSU       & CPSG\\
Transfer (engineering)   & SPJU       & SPJG ($\times 5$)\\
Transfer (synthesis)     & CPJU       & CPJG ($\times 5$)\\
Creation / frontier      & ---        & WRQ ($\times 5$), ORQ ($\times 5$)\\
\bottomrule
\end{tabularx}
\end{center}

\newpage
%=============================================================
\subsection{Undergraduate Track Artifacts}
\label{sec:ug-artifacts}
%=============================================================

\begin{tcolorbox}[trackheaderUG]
UNDERGRADUATE TRACK \ \ (Tiers 1--3: HS, BegUG, AdvUG)
\end{tcolorbox}

\bigskip

%--- SCQU ----------------------------------------------------
\subsubsection{SCQU --- Simple Concept Questions (Undergrad)}

\begin{tcolorbox}[ugconceptbox, title={SCQU --- Simple Concept Questions, Undergraduate \hfill\normalfont 10 questions}]

\textbf{Purpose:} Fast checks that a Tier~1--3 student has the \emph{right mental model},
not just algebra. Pitched at intuition and working-language levels.

\medskip\textbf{Format:}
\begin{itemize}[leftmargin=1.5em,itemsep=2pt]
  \item 10 questions --- mix of multiple-choice (4), true/false + one-sentence justify (4),
        and 1--2 sentence short answer (2)
  \item 15--20 minutes in class or as pre-lecture warmup
\end{itemize}

\medskip\textbf{Tier targeting:}
\begin{center}
\begin{tabular}{lp{9cm}}
\toprule
HS     & Qualitative only; language avoids formalism; diagrams welcome\\
BegUG  & Qualitative + brief algebraic check (e.g.\ verify $|c_1|^2+|c_2|^2=1$)\\
AdvUG  & Include ``identify the correct statement'' from near-equivalent options\\
\bottomrule
\end{tabular}
\end{center}

\medskip\textbf{Example prompts (L01 style):}
\begin{itemize}[leftmargin=1.5em,itemsep=2pt]
  \item Why is $|\psi\rangle$ not ``a column vector'' until a basis is chosen?
  \item What is physically meaningful: $\langle a|\psi\rangle$ or $|\langle a|\psi\rangle|^2$?
  \item True/False: multiplying $|\psi\rangle$ by $e^{i\theta}$ changes the probability of
        any outcome. (Justify.)
  \item In what sense is ``superposition'' basis-dependent?
\end{itemize}
\end{tcolorbox}

\begin{tcolorbox}[rubricbox, title={SCQU Rubric}]
\begin{tabular}{ll}
  Correct answer                & 1.5 pts each\\
  Brief justification consistent & 0.5 pts each\\
\end{tabular}\\[4pt]
\textbf{Total: 20 points} (10 questions $\times$ 2 pts)
\end{tcolorbox}

\bigskip
%--- CCQU ----------------------------------------------------
\subsubsection{CCQU --- Complex Concept Questions (Undergrad)}

\begin{tcolorbox}[ugconceptbox, title={CCQU --- Complex Concept Questions, Undergraduate \hfill\normalfont 10 questions}]

\textbf{Purpose:} Test \emph{structural understanding} at the UG level: not just ``what''
but ``why the formalism is consistent.'' Requires short arguments, not full proofs.

\medskip\textbf{Format:}
\begin{itemize}[leftmargin=1.5em,itemsep=2pt]
  \item 10 questions --- mix of ``explain why'' (4), ``what would change if\ldots'' (3),
        and ``identify the flaw in this argument'' (3)
  \item 25--35 minutes; homework or tutorial context
\end{itemize}

\medskip\textbf{Tier targeting:}
\begin{center}
\begin{tabular}{lp{9cm}}
\toprule
BegUG  & Explain consequences of definitions; trace through an argument\\
AdvUG  & Identify logical gaps; construct a counterexample; compare two approaches\\
\bottomrule
\end{tabular}
\end{center}

\medskip\textbf{Example prompts (L01 style):}
\begin{itemize}[leftmargin=1.5em,itemsep=2pt]
  \item Why is the global phase of a quantum state unobservable, but the \emph{relative}
        phase between components is not? Give a concrete example.
  \item A student claims: ``$|\psi\rangle = \tfrac{1}{\sqrt{2}}(|0\rangle+|1\rangle)$
        and $|\phi\rangle = \tfrac{1}{\sqrt{2}}(|0\rangle-|1\rangle)$ describe the
        same physical state.'' What is wrong with this claim?
  \item If Hilbert space were defined over the real numbers instead of complex numbers,
        which aspect of quantum mechanics would immediately fail? Justify briefly.
\end{itemize}
\end{tcolorbox}

\begin{tcolorbox}[rubricbox, title={CCQU Rubric}]
\begin{tabular}{ll}
  Correct identification of the key point & 1.0 pt each\\
  Coherent short argument or example       & 1.0 pt each\\
\end{tabular}\\[4pt]
\textbf{Total: 20 points} (10 questions $\times$ 2 pts)
\end{tcolorbox}

\bigskip
%--- SPSU ----------------------------------------------------
\subsubsection{SPSU --- Simple Problem Set (Undergrad)}

\begin{tcolorbox}[ugprobbox, title={SPSU --- Simple Problem Set, Undergraduate \hfill\normalfont 10 problems}]

\textbf{Purpose:} Build \emph{procedural fluency}: compute probabilities, normalise
states, insert identity, perform basic basis changes.

\medskip\textbf{Format:} 10 problems; 1--5 steps each; minimal proofs.

\medskip\textbf{Internal structure:}
\begin{itemize}[leftmargin=1.5em,itemsep=2pt]
  \item \textbf{Part A --- mechanics} (4 problems): normalisation, components, overlaps
  \item \textbf{Part B --- interpretation} (3 problems): what does the computed number mean?
  \item \textbf{Part C --- translation} (3 problems):
        ket $\leftrightarrow$ coordinates $\leftrightarrow$ wavefunction preview
\end{itemize}

\medskip\textbf{Tier scaling:}
\begin{center}
\begin{tabular}{lp{8.5cm}}
\toprule
HS    & Numeric complex arithmetic + amplitude-to-probability conversion\\
BegUG & Normalisation + components $c_n=\langle n|\psi\rangle$ + probability sums\\
AdvUG & Include representation-change: verify result independent of basis choice\\
\bottomrule
\end{tabular}
\end{center}
\end{tcolorbox}

\begin{tcolorbox}[rubricbox, title={SPSU Rubric (per problem)}]
\begin{tabular}{ll}
  Correct result           & 60\%\\
  Correct method/notation  & 30\%\\
  Interpretation sentence  & 10\%\\
\end{tabular}
\end{tcolorbox}

\bigskip
%--- CPSU ----------------------------------------------------
\subsubsection{CPSU --- Complex Problem Set (Undergrad)}

\begin{tcolorbox}[ugprobbox, title={CPSU --- Complex Problem Set, Undergraduate \hfill\normalfont 5 problems}]

\textbf{Purpose:} Assess \emph{structural mastery} at UG level: proofs, derivations,
and multi-step synthesis accessible to Tier~3 students.

\medskip\textbf{Format:} 5 multi-part problems (a)(b)(c)\ldots; 3--6 hours total.

\medskip\textbf{Canonical themes:}
\begin{itemize}[leftmargin=1.5em,itemsep=2pt]
  \item Prove Cauchy--Schwarz and/or triangle inequality
  \item Prove phase invariance and ray equivalence (with guidance)
  \item Show unitaries preserve probabilities and inner products
  \item Derive completeness consequences (Parseval identity)
  \item Construct a worked proof of the spectral theorem for a $2\times 2$ Hermitian matrix
\end{itemize}

\medskip\textbf{Tier targeting:} Primarily AdvUG (Tier~3); parts (a) accessible to BegUG.
\end{tcolorbox}

\begin{tcolorbox}[rubricbox, title={CPSU Rubric (per problem)}]
\begin{tabular}{ll}
  Correct statement of assumptions & 15\%\\
  Logical proof structure           & 45\%\\
  Correct algebra/steps             & 25\%\\
  Insight / interpretation          & 15\%\\
\end{tabular}
\end{tcolorbox}

\bigskip
%--- SPJU ----------------------------------------------------
\subsubsection{SPJU --- Simple Project (Undergrad)}

\begin{tcolorbox}[ugprobbox, title={SPJU --- Simple Project, Undergraduate \hfill\normalfont 1 project}]

\textbf{Purpose:} Hands-on transfer: ``I can build something that uses the formalism.''

\medskip\textbf{Scope:} 2--6 hours; structured, scaffolded; works for mixed Tier~1--3 cohort.

\medskip\textbf{Deliverable:} Short report (1--3 pages) \emph{and} a small artifact
(code, worksheet, or visualisation).

\medskip\textbf{Example --- Basis Expansion Calculator (L01):}
\begin{itemize}[leftmargin=1.5em,itemsep=2pt]
  \item \textit{Inputs:} coefficients $c$ in basis $\{|n\rangle\}$; unitary $U$
  \item \textit{Outputs:} new coefficients $c'=Uc$
  \item \textit{Required checks:} normalisation preserved; probabilities sum to 1;
        global phase invariance verified
  \item \textit{Deliverables:} code + README; one worked example; short explanation
        of what changed and what did not
\end{itemize}
\end{tcolorbox}

\begin{tcolorbox}[rubricbox, title={SPJU Rubric}]
\begin{tabular}{ll}
  Correctness                    & 40\%\\
  Verification tests included    & 25\%\\
  Clarity of explanation         & 25\%\\
  Reproducibility (instructions) & 10\%\\
\end{tabular}
\end{tcolorbox}

\bigskip
%--- CPJU ----------------------------------------------------
\subsubsection{CPJU --- Complex Project (Undergrad)}

\begin{tcolorbox}[ugprobbox, title={CPJU --- Complex Project, Undergraduate \hfill\normalfont 1 project}]

\textbf{Purpose:} Synthesis and conceptual transfer: ``I can compare frameworks and
defend choices.''

\medskip\textbf{Scope:} 8--15 hours; open-ended but guided; mini-essay or mini-tool
with reflection.

\medskip\textbf{Deliverable:} 5--8 page report + appendix with computations +
optional code.

\medskip\textbf{Example --- Compare Representations (L01):}
\begin{enumerate}[leftmargin=2em,itemsep=2pt]
  \item Represent a chosen state as: (i) abstract ket, (ii) coordinates in basis A,
        (iii) coordinates in basis B, (iv) wavefunction preview $\psi(x)=\langle x|\psi\rangle$
  \item Analyse: what is invariant (norm, probabilities)? What changes (coordinates)?
\end{enumerate}
\end{tcolorbox}

\begin{tcolorbox}[rubricbox, title={CPJU Rubric}]
\begin{tabular}{ll}
  Technical correctness   & 30\%\\
  Conceptual synthesis    & 35\%\\
  Depth \& structure      & 25\%\\
  Originality / insight   & 10\%\\
\end{tabular}
\end{tcolorbox}

\newpage
%=============================================================
\subsection{Graduate Track Artifacts}
\label{sec:grad-artifacts}
%=============================================================

\begin{tcolorbox}[trackheaderGRAD]
GRADUATE TRACK \ \ (Tiers 4--5: MSc, PhD)
\end{tcolorbox}

\bigskip

%--- SCQG ----------------------------------------------------
\subsubsection{SCQG --- Simple Concept Questions (Grad)}

\begin{tcolorbox}[gradconceptbox, title={SCQG --- Simple Concept Questions, Graduate \hfill\normalfont 10 questions}]

\textbf{Purpose:} Rapid check that a Tier~4--5 student can \emph{state and navigate}
the rigorous formalism without hesitation.

\medskip\textbf{Format:} 10 questions; mix of short statement (4), true/false with proof
sketch (3), and ``identify the precise error'' (3); 20--25 minutes.

\medskip\textbf{Tier targeting:}
\begin{center}
\begin{tabular}{lp{9cm}}
\toprule
MSc & Conceptual statements about rays, operators, measurement; standard results\\
PhD & Include precise domain/boundedness conditions; statement of named theorems\\
\bottomrule
\end{tabular}
\end{center}

\medskip\textbf{Example prompts (L01 style):}
\begin{itemize}[leftmargin=1.5em,itemsep=2pt]
  \item State precisely what it means for a Hilbert space to be separable.
  \item True/False (with proof sketch): every Hermitian operator on $\mathcal{H}$ is
        self-adjoint.
  \item What is a ray in $\mathcal{H}$, and why are rays the correct mathematical objects
        for pure quantum states?
  \item Give the exact statement of Wigner's theorem (including the $U(1)$ ambiguity).
\end{itemize}
\end{tcolorbox}

\begin{tcolorbox}[rubricbox, title={SCQG Rubric}]
\begin{tabular}{ll}
  Precise correct statement       & 1.5 pts each\\
  Brief justification or proof sketch & 0.5 pts each\\
\end{tabular}\\[4pt]
\textbf{Total: 20 points}
\end{tcolorbox}

\bigskip
%--- CCQG ----------------------------------------------------
\subsubsection{CCQG --- Complex Concept Questions (Grad)}

\begin{tcolorbox}[gradconceptbox, title={CCQG --- Complex Concept Questions, Graduate \hfill\normalfont 10 questions}]

\textbf{Purpose:} Test \emph{deep structural understanding} at graduate level: connecting
formalism to mathematical foundations, identifying subtle distinctions, and engaging with
open or contested points.

\medskip\textbf{Format:} 10 questions; mix of ``distinguish and explain'' (4),
``what fails without assumption X'' (3), and ``connect these two results'' (3);
35--50 minutes.

\medskip\textbf{Example prompts (L01 style):}
\begin{itemize}[leftmargin=1.5em,itemsep=2pt]
  \item Distinguish Hermitian from self-adjoint operators. Give an example of an operator
        that is one but not the other.
  \item Why does the Fubini--Study metric on $\mathbb{P}(\mathcal{H})$ encode transition
        probabilities? Show the connection explicitly.
  \item What additional structure, beyond a Hilbert space, is needed to define the
        Gel'fand triple? Why is this structure physically necessary?
  \item Wigner's theorem guarantees every symmetry lifts to a unitary or antiunitary
        operator. What happens if we drop the assumption that the map preserves
        $|\langle\phi|\psi\rangle|^2$? Is there still a useful structure theorem?
\end{itemize}
\end{tcolorbox}

\begin{tcolorbox}[rubricbox, title={CCQG Rubric}]
\begin{tabular}{ll}
  Correct identification of key distinction/connection & 1.0 pt each\\
  Rigorous short argument or counterexample            & 1.0 pt each\\
\end{tabular}\\[4pt]
\textbf{Total: 20 points}
\end{tcolorbox}

\bigskip
%--- SPSG ----------------------------------------------------
\subsubsection{SPSG --- Simple Problem Set (Grad)}

\begin{tcolorbox}[gradprobbox, title={SPSG --- Simple Problem Set, Graduate \hfill\normalfont 10 problems}]

\textbf{Purpose:} Build \emph{fluency with the rigorous formalism}: operator domains,
spectral theory computations, density matrices, and representation-independence checks.

\medskip\textbf{Format:} 10 problems; each 2--8 steps; representation-independence
always verified.

\medskip\textbf{Internal structure:}
\begin{itemize}[leftmargin=1.5em,itemsep=2pt]
  \item \textbf{Part A --- operator mechanics} (4 problems): adjoint, spectrum,
        domain checks, expectation values in density-matrix form
  \item \textbf{Part B --- rigorous translation} (3 problems): continuous bases,
        delta distributions, Fourier connection
  \item \textbf{Part C --- symmetry \& invariance} (3 problems): commutator algebra,
        unitary equivalence, Wigner theorem applications
\end{itemize}

\medskip\textbf{Tier scaling:}
\begin{center}
\begin{tabular}{lp{8.5cm}}
\toprule
MSc & Work in density-matrix formalism; verify trace-class and positivity conditions\\
PhD & Additionally check domain conditions; apply functional calculus\\
\bottomrule
\end{tabular}
\end{center}
\end{tcolorbox}

\begin{tcolorbox}[rubricbox, title={SPSG Rubric (per problem)}]
\begin{tabular}{ll}
  Correct result with domain/condition check & 50\%\\
  Correct method and rigorous notation       & 35\%\\
  Interpretation or physical meaning         & 15\%\\
\end{tabular}
\end{tcolorbox}

\bigskip
%--- CPSG ----------------------------------------------------
\subsubsection{CPSG --- Complex Problem Set (Grad)}

\begin{tcolorbox}[gradprobbox, title={CPSG --- Complex Problem Set, Graduate \hfill\normalfont 5 problems}]

\textbf{Purpose:} Assess \emph{graduate-level structural mastery}: complete rigorous
proofs, functional analysis arguments, and graduate-level derivations.

\medskip\textbf{Format:} 5 multi-part problems; 5--10 hours total; no guidance scaffolding.

\medskip\textbf{Canonical themes:}
\begin{itemize}[leftmargin=1.5em,itemsep=2pt]
  \item Prove the nuclear spectral theorem in rigged Hilbert space context
  \item Prove Stone's theorem (statement + sketch of key steps)
  \item Prove that $\hat{p}=-i\hbar d/dx$ is self-adjoint on $H^1(\mathbb{R})$
        but only symmetric on $L^2([0,\infty))$
  \item Prove the Riesz representation theorem and connect to the bra as linear functional
  \item Derive the Fubini--Study metric on $\mathbb{CP}^n$ from the Hilbert space
        inner product
\end{itemize}
\end{tcolorbox}

\begin{tcolorbox}[rubricbox, title={CPSG Rubric (per problem)}]
\begin{tabular}{ll}
  Precise statement of hypotheses         & 15\%\\
  Logical proof structure                 & 40\%\\
  Correct analysis / functional arguments & 30\%\\
  Contextual insight / connections        & 15\%\\
\end{tabular}
\end{tcolorbox}

\bigskip
%--- SPJG ----------------------------------------------------
\subsubsection{SPJG --- Simple Projects (Grad)}

\begin{tcolorbox}[gradprobbox, title={SPJG --- Simple Projects, Graduate \hfill\normalfont 5 projects}]

\textbf{Purpose:} Structured transfer at graduate level: ``I can implement and test a
rigorous idea.'' Five independent projects, each with a clear deliverable.

\medskip\textbf{Scope per project:} 3--8 hours; structured, with defined inputs and outputs.

\medskip\textbf{Deliverable per project:} Technical note (2--4 pages) + code or
formal derivation document.

\medskip\textbf{Typical SPJG project types (L01 style):}
\begin{enumerate}[leftmargin=2em,itemsep=4pt]
  \item \textbf{Operator domain explorer:} For $\hat{p}$ on various domains, classify
        Hermitian/symmetric/self-adjoint; tabulate deficiency indices.
  \item \textbf{Density matrix diagnostics:} Given $4\times 4$ matrices, test
        $\hat{\rho}\ge 0$, $\operatorname{Tr}\hat{\rho}=1$, purity; classify mixed/pure/entangled.
  \item \textbf{Wigner function sampler:} Compute and plot the Wigner quasi-probability
        for a set of states; identify negativity signatures.
  \item \textbf{Spectral measure calculator:} For a simple Hamiltonian, construct the
        spectral projection $E_\lambda$ and verify the spectral theorem numerically.
  \item \textbf{Ray geometry visualiser:} Implement and visualise the Fubini--Study
        metric for $\mathbb{CP}^1$ (Bloch sphere); compute geodesic distances between
        given states.
\end{enumerate}
\end{tcolorbox}

\begin{tcolorbox}[rubricbox, title={SPJG Rubric (per project)}]
\begin{tabular}{ll}
  Mathematical correctness               & 40\%\\
  Rigour of conditions/checks            & 25\%\\
  Clarity of technical note              & 25\%\\
  Reproducibility                        & 10\%\\
\end{tabular}
\end{tcolorbox}

\bigskip
%--- CPJG ----------------------------------------------------
\subsubsection{CPJG --- Complex Projects (Grad)}

\begin{tcolorbox}[gradprobbox, title={CPJG --- Complex Projects, Graduate \hfill\normalfont 5 projects}]

\textbf{Purpose:} Open synthesis at graduate level: ``I can connect frameworks, identify
limitations, and produce research-quality output.''
Five independent projects, each requiring significant independent work.

\medskip\textbf{Scope per project:} 10--20 hours; open-ended but guided by a question.

\medskip\textbf{Deliverable per project:} 6--10 page technical report + appendix.

\medskip\textbf{Typical CPJG project types (L01 style):}
\begin{enumerate}[leftmargin=2em,itemsep=4pt]
  \item \textbf{Real vs.\ complex Hilbert space:} Identify which quantum phenomena
        fail in a real Hilbert space; connect to Wigner's theorem over $\mathbb{R}$.
  \item \textbf{Self-adjoint extensions:} For $\hat{p}$ on $[0,1]$ with varied boundary
        conditions, classify all self-adjoint extensions; compute their spectra.
  \item \textbf{Rigged Hilbert space construction:} Build explicitly the Gel'fand triple
        $\Phi\subset L^2(\mathbb{R})\subset\Phi'$ for the harmonic oscillator; identify
        the generalised eigenvectors of $\hat{x}$.
  \item \textbf{Projective space and entanglement:} Characterise the image of the
        Segre embedding $\mathbb{CP}^1\times\mathbb{CP}^1\hookrightarrow\mathbb{CP}^3$;
        interpret the complement as entangled states.
  \item \textbf{Measurement model analysis:} Compare von Neumann, L\"uders, and
        POVM update rules; identify conditions under which they agree and diverge.
\end{enumerate}
\end{tcolorbox}

\begin{tcolorbox}[rubricbox, title={CPJG Rubric (per project)}]
\begin{tabular}{ll}
  Technical rigour and correctness  & 30\%\\
  Conceptual synthesis              & 30\%\\
  Depth, structure, and exposition  & 25\%\\
  Originality / critical insight    & 15\%\\
\end{tabular}
\end{tcolorbox}

\newpage
%=============================================================
\subsection{Research Track Artifacts}
\label{sec:research-artifacts}
%=============================================================

\begin{tcolorbox}[trackheaderRES]
RESEARCH TRACK \ \ (Tiers 4--5: MSc, PhD) \ \ shared across both graduate tracks
\end{tcolorbox}

\bigskip

%--- WRQ ----------------------------------------------------
\subsubsection{WRQ --- Well-Defined Research Questions ($\times$5)}

\begin{tcolorbox}[researchbox, title={WRQ --- Well-Defined Research Questions \hfill\normalfont 5 questions}]

\textbf{Purpose:} Train ``research-like thinking'' while remaining \emph{bounded}:
clear scope, clear deliverable, answerable within a defined time.

\medskip\textbf{Scope per question:} 5--10 hours; answerable with literature +
calculation + interpretation.

\medskip\textbf{Deliverable per question --- 2--6 page memo:}
\begin{itemize}[leftmargin=1.5em,itemsep=2pt]
  \item Question statement
  \item Method (how you will answer it)
  \item Evidence (citations + computations)
  \item Conclusion
  \item 5--10 references
\end{itemize}

\medskip\textbf{Five example WRQ prompts (Module I.2, bra-ket context):}
\begin{enumerate}[leftmargin=2em,itemsep=4pt]
  \item ``How did Born's probabilistic interpretation change the meaning of the
        wavefunction between 1926--1927? Trace the primary sources.''
  \item ``What exactly does `state is a ray' buy us operationally? Provide
        examples where vectors fail but rays succeed.''
  \item ``When and why did von Neumann introduce the density operator? What problem
        was it designed to solve that pure states cannot address?''
  \item ``What is the minimal set of axioms on an inner product space that forces
        the Born rule to take the form $P=|\langle a|\psi\rangle|^2$? Compare
        Gleason's theorem with the Dirac--von Neumann axioms.''
  \item ``How does the choice of complex vs.\ real Hilbert space affect the
        structure of pure-state spaces? What operational consequence follows?''
\end{enumerate}
\end{tcolorbox}

\begin{tcolorbox}[rubricbox, title={WRQ Rubric (per question)}]
\begin{tabular}{ll}
  Clear scope and method    & 25\%\\
  Evidence quality          & 25\%\\
  Correct technical content & 25\%\\
  Quality of argument       & 25\%\\
\end{tabular}
\end{tcolorbox}

\bigskip
%--- ORQ ----------------------------------------------------
\subsubsection{ORQ --- Open-Ended Research Questions ($\times$5)}

\begin{tcolorbox}[researchbox, title={ORQ --- Open-Ended Research Questions \hfill\normalfont 5 questions}]

\textbf{Purpose:} Expose the frontier. No single correct answer; evaluates
\emph{quality of reasoning}, not correctness of conclusion.

\medskip\textbf{Scope per question:} 10--20 hours (or ongoing); may include proposing
a new definition, comparing frameworks, or designing a thought experiment.

\medskip\textbf{Deliverable per question --- 4--10 page proposal-style note:}
\begin{itemize}[leftmargin=1.5em,itemsep=2pt]
  \item Motivation
  \item What is known (brief review)
  \item What you propose or conjecture
  \item How you would test or argue it
  \item Risks and limitations
\end{itemize}

\medskip\textbf{Five example ORQ prompts (Module I.2, bra-ket context):}
\begin{enumerate}[leftmargin=2em,itemsep=4pt]
  \item ``What minimal operational axioms force complex Hilbert space rather than
        real or quaternionic? Is the complex structure fundamental or emergent?''
  \item ``Can you formulate measurement and state update without projective collapse
        but still reproduce all observed statistics? What would such a framework
        look like mathematically?''
  \item ``The Gel'fand triple extends the Hilbert space to accommodate
        distributions. Can you propose an alternative foundation that avoids
        distributions entirely while retaining the continuous spectrum?''
  \item ``Wigner's theorem guarantees symmetries are unitary or antiunitary.
        Are there physically motivated generalisations (e.g.\ to mixed states, open
        systems, or approximate symmetries) that require a different structural
        result?''
  \item ``Is separability of the Hilbert space a physical assumption or merely a
        mathematical convenience? Design an operational test that would distinguish
        a separable from a non-separable state space in principle.''
\end{enumerate}
\end{tcolorbox}

\begin{tcolorbox}[rubricbox, title={ORQ Rubric (per question)}]
\begin{tabular}{ll}
  Clarity and framing of the question  & 25\%\\
  Depth of engagement with literature  & 25\%\\
  Novelty and plausibility of proposal & 25\%\\
  Awareness of limitations / open ends & 25\%\\
\end{tabular}
\end{tcolorbox}

\newpage
%=============================================================
\subsection{Exercise Packaging Standard}
%=============================================================

\instructornote{Every individual problem in CPSU, CPSG, or any project type should
be packaged in the four-layer format below. This supports self-study, tiered hints
delivery, and automated repository assembly.}

\begin{tcolorbox}[goldbox]
\textbf{Four-layer exercise packet} (one file-set per problem):

\medskip
\begin{tabular}{lp{9.5cm}}
\toprule
\textbf{Layer} & \textbf{Content}\\
\midrule
1. Problem Statement    & Clean statement; no hints; the only file students receive initially\\
2. Suggestions / Hints  & Scaffolded prompts; released after first attempt\\
3. Advanced Suggestions & Structural hints, alternative approaches, cross-references\\
4. Detailed Solution    & Full derivation + commentary + common pitfall warnings\\
\bottomrule
\end{tabular}

\medskip\noindent\textbf{File naming convention:}\\
\texttt{ex\#\#\_problem.tex} \quad
\texttt{ex\#\#\_hints.tex} \quad
\texttt{ex\#\#\_adv\_hints.tex} \quad
\texttt{ex\#\#\_solution.tex}
\end{tcolorbox}

%=============================================================
\subsection{Per-Lecture Minimum Bundle}
%=============================================================

\begin{tcolorbox}[goldbox]
\textbf{Complete artifact bundle for one 90-minute lecture:}

\medskip
\begin{center}
\begin{tabular}{llrr}
\toprule
\textbf{Code} & \textbf{Name} & \textbf{Track} & \textbf{Count}\\
\midrule
SCQU & Simple Concept Questions --- UG   & UG       & 10 questions\\
CCQU & Complex Concept Questions --- UG  & UG       & 10 questions\\
SPSU & Simple Problem Set --- UG         & UG       & 10 problems\\
CPSU & Complex Problem Set --- UG        & UG       & 5 problems\\
SPJU & Simple Project --- UG             & UG       & 1 project\\
CPJU & Complex Project --- UG            & UG       & 1 project\\
\midrule
SCQG & Simple Concept Questions --- Grad & Grad     & 10 questions\\
CCQG & Complex Concept Questions --- Grad& Grad     & 10 questions\\
SPSG & Simple Problem Set --- Grad       & Grad     & 10 problems\\
CPSG & Complex Problem Set --- Grad      & Grad     & 5 problems\\
SPJG & Simple Projects --- Grad         & Grad     & 5 projects\\
CPJG & Complex Projects --- Grad        & Grad     & 5 projects\\
\midrule
WRQ  & Well-Defined Research Questions   & Research & 5 questions\\
ORQ  & Open-Ended Research Questions     & Research & 5 questions\\
\midrule
 & \textbf{Total per lecture}           &          & \textbf{92 items}\\
\bottomrule
\end{tabular}
\end{center}
\end{tcolorbox}

%=============================================================
\subsection{Suggested Assessment Weights}
%=============================================================

\begin{center}
\begin{tabular}{llrp{5.5cm}}
\toprule
\textbf{Track} & \textbf{Artifact group} & \textbf{Weight} & \textbf{Rationale}\\
\midrule
\multirow{4}{*}{\textcolor{ugblue}{UG}}
  & SCQU + CCQU (concept)  & 15\% & Mental models and structural understanding\\
  & SPSU (fluency)         & 25\% & Core procedural skills; highest volume\\
  & CPSU (mastery)         & 30\% & Proof and derivation ability\\
  & SPJU + CPJU (transfer) & 30\% & Application and synthesis\\
\midrule
\multirow{4}{*}{\textcolor{gradred}{Grad}}
  & SCQG + CCQG (concept)  & 15\% & Rigorous conceptual command\\
  & SPSG (fluency)         & 20\% & Formal computation fluency\\
  & CPSG (mastery)         & 30\% & Graduate proof standards\\
  & SPJG + CPJG (transfer) & 35\% & Research-quality synthesis\\
\midrule
\multirow{2}{*}{\textcolor{respath}{Research}}
  & WRQ ($\times$5)        & 50\% & Scoped research capacity\\
  & ORQ ($\times$5)        & 50\% & Frontier reasoning\\
\bottomrule
\end{tabular}
\end{center}

\newpage
%=============================================================
\section*{Appendix --- Quick Reference Card}
\addcontentsline{toc}{section}{Appendix --- Quick Reference Card}
%=============================================================

\begin{center}
{\large\textbf{Tier $\times$ Artifact Eligibility Grid}}\\[8pt]
\small
\begin{tabular}{l|cc|cc|cc|cc|cc|cc|cc}
\toprule
 & \multicolumn{2}{c|}{\textbf{Concept Q.}}
 & \multicolumn{2}{c|}{\textbf{Problem Set}}
 & \multicolumn{2}{c|}{\textbf{Projects}}
 & \multicolumn{2}{c|}{\textbf{Concept Q.}}
 & \multicolumn{2}{c|}{\textbf{Problem Set}}
 & \multicolumn{2}{c|}{\textbf{Projects}}
 & \multicolumn{2}{c}{\textbf{Research}}\\
 & S & C & S & C & S & C & S & C & S & C & S & C & W & O\\
 & C & C & P & P & P & P & C & C & P & P & J & J & R & R\\
 & Q & Q & S & S & J & J & Q & Q & S & S & G & G & Q & Q\\
 & U & U & U & U & U & U & G & G & G & G & & & &\\
\midrule
HS     & \checkmark &            & \checkmark &            & \checkmark &            &            &            &            &            &            &            &            &\\
BegUG  & \checkmark & \checkmark & \checkmark &            & \checkmark & opt.       &            &            &            &            &            &            &            &\\
AdvUG  & \checkmark & \checkmark & \checkmark & \checkmark & \checkmark & \checkmark &            &            &            &            &            &            &            &\\
MSc    &            &            &            &            &            &            & \checkmark & \checkmark & \checkmark & \checkmark & \checkmark & \checkmark & \checkmark & \checkmark\\
PhD    &            &            &            &            &            &            & \checkmark & \checkmark & \checkmark & \checkmark & \checkmark & \checkmark & \checkmark & \checkmark\\
\bottomrule
\end{tabular}
\end{center}

\bigskip
\begin{center}
{\large\textbf{Artifact Duration and Output Summary}}\\[8pt]
\begin{tabular}{llllr}
\toprule
\textbf{Code} & \textbf{Format} & \textbf{Est.\ time} & \textbf{Primary output} & \textbf{Pts}\\
\midrule
SCQU & 10 questions & 15--20 min  & Answers + brief justifications & 20\\
CCQU & 10 questions & 25--35 min  & Short arguments + examples     & 20\\
SPSU & 10 problems  & 2--4 hrs    & Worked solutions               & 100\\
CPSU & 5 problems   & 3--6 hrs    & Proofs + derivations           & 50\\
SPJU & 1 project    & 2--6 hrs    & Code/worksheet + 1--3 pg report& 100\\
CPJU & 1 project    & 8--15 hrs   & 5--8 pg report + appendix      & 100\\
\midrule
SCQG & 10 questions & 20--25 min  & Precise statements + sketches  & 20\\
CCQG & 10 questions & 35--50 min  & Rigorous arguments             & 20\\
SPSG & 10 problems  & 3--6 hrs    & Formal solutions + checks      & 100\\
CPSG & 5 problems   & 5--10 hrs   & Graduate-level proofs          & 50\\
SPJG & 5 projects   & 3--8 hrs ea.& Tech notes + code              & 5$\times$100\\
CPJG & 5 projects   & 10--20 hrs ea.& Full technical reports       & 5$\times$100\\
\midrule
WRQ  & 5 memos      & 5--10 hrs ea. & 2--6 pg research memo       & 5$\times$100\\
ORQ  & 5 proposals  & 10--20 hrs ea.& 4--10 pg proposal note      & 5$\times$100\\
\bottomrule
\end{tabular}
\end{center}

\newpage
%#############################################################
\section{Tiered Bibliography --- Structure \& Design Principles}
\label{sec:bibdesign}
%#############################################################

\begin{tcolorbox}[goldbox]
\textbf{This section explains how the companion bibliography document is built and why.}
The full reading lists --- one per learning tier, plus a complete alphabetical catalogue
--- are published as a separate document: \textbf{\textit{QM Programme: Tiered Bibliography}}.
This section gives instructors and students the conceptual framework needed to
\emph{use}, \emph{extend}, and \emph{adapt} that bibliography for other modules and disciplines.
\end{tcolorbox}

%-------------------------------------------------------------
\subsection{Purpose of a Tiered Bibliography}
%-------------------------------------------------------------

A flat reading list attached to a syllabus has a fundamental problem: it collapses
five very different audiences into one list. A Tier~1 student who opens Reed \& Simon's
\textit{Methods of Modern Mathematical Physics} will be lost; a PhD student assigned
Feynman's \textit{QED} may feel patronised. The tiered bibliography solves this by
providing \textbf{one list per tier}, each answering a single question:
\begin{quote}\itshape
  What should I read \emph{next}, at \emph{my} level, to go deeper into the material
  of this lecture?
\end{quote}

\noindent The tier structure mirrors exactly the five-tier pedagogy described in
Section~\ref{sec:tiers} of this document, and the reading strategy at each tier
aligns with the artifact types described in Section~\ref{sec:artifacts}.

%-------------------------------------------------------------
\subsection{Four Categories of Material}
%-------------------------------------------------------------

Every entry in the bibliography belongs to exactly one of four categories,
marked by a tag in bold square brackets:

\begin{center}
\begin{tabular}{lp{10cm}}
\toprule
\textbf{Tag} & \textbf{Category and description}\\
\midrule
\textbf{[N]} & \textit{Popular / narrative.} Builds intuition and historical context;
               no formal prerequisites. Assigned at Tier~1; useful context at Tier~2.\\[4pt]
\textbf{[T]} & \textit{Pedagogical textbook.} Structured course reading with exercises;
               the backbone of Tiers~1--4. These are the books students read
               chapter-by-chapter alongside lectures.\\[4pt]
\textbf{[P]} & \textit{Primary research paper.} Original source; accessed at Tier~3
               (abstract + introduction), read fully from Tier~4; mandatory at Tier~5.
               The companion bibliography identifies 14 primary papers spanning
               quantum mechanics and nonlinear dynamics.\\[4pt]
\textbf{[M]} & \textit{Research monograph or review article.} Graduate-level and
               research-level consolidation; Tiers~4--5 and the Research track.
               Monographs are read selectively (relevant chapters); reviews are
               read in full as literature surveys.\\
\bottomrule
\end{tabular}
\end{center}

%-------------------------------------------------------------
\subsection{Annotation System}
%-------------------------------------------------------------

Each entry in the tiered bibliography carries three elements beyond the citation:

\medskip
\begin{center}
\begin{tabular}{lp{10.5cm}}
\toprule
\textbf{Element} & \textbf{Meaning and use}\\
\midrule
Star rating & \textbf{\large$\bigstar\bigstar\bigstar$} Essential --- read in full
             \quad\textbf{\large$\bigstar\bigstar$} Strongly recommended --- read
             relevant chapters
             \quad\textbf{\large$\bigstar$} Useful supplement --- consult as needed\\[6pt]
Type tag    & [N] [T] [P] [M] --- see categories above\\[6pt]
Annotation  & 2--4 sentences linking the entry to specific programme modules and
             lecture codes (e.g.\ ``Chapters~1--3 support Module~I.2 Lectures L01--L05'').
             This is what distinguishes a tiered programme bibliography from a generic
             reading list: every entry is \emph{placed} in the curriculum.\\
\bottomrule
\end{tabular}
\end{center}

%-------------------------------------------------------------
\subsection{How the Tiers Map to Reading Strategies}
%-------------------------------------------------------------

\begin{center}
\begin{tabular}{clp{3.8cm}p{5cm}}
\toprule
\textbf{Tier} & \textbf{Label} & \textbf{Primary material type} & \textbf{Reading strategy}\\
\midrule
1 & HS       & [N] popular books     & Read cover-to-cover before the course\\
2 & BegUG    & [T] textbooks         & One chapter per lecture; do all exercises\\
3 & AdvUG    & [T] + [P] first reads & Textbook main; read abstracts/intros of primary papers\\
4 & MSc      & [P] + [M] selected    & Primary papers as main reading; textbooks as reference\\
5 & PhD      & [M] + [P] all         & Monographs + full primary literature + arXiv\\
\bottomrule
\end{tabular}
\end{center}

\medskip
The tiers are \textbf{cumulative}: Tier~$k$ reading is assumed as background for
Tier~$k{+}1$. Instructors assigning material to a mixed cohort can use the tier badges
in the bibliography to route students automatically.

%-------------------------------------------------------------
\subsection{Structure of the Companion Bibliography Document}
%-------------------------------------------------------------

The companion document \textbf{\textit{QM Programme: Tiered Bibliography}} is organised
as follows:

\begin{center}
\begin{tabular}{lp{10cm}}
\toprule
\textbf{Part} & \textbf{Content}\\
\midrule
Introduction & How-to-use guide; annotation key; tier--strategy table\\
Section~1 (T1) & Popular science \& first introductions --- QM + Nonlinear Dynamics\\
Section~2 (T2) & Core undergraduate textbooks --- QM + Mathematics support + QChem\\
Section~3 (T3) & Rigorous undergraduate \& bridge texts --- QM + primary papers (first encounter)\\
Section~4 (T4) & Graduate textbooks, many-body theory, quantum chemistry + key papers\\
Section~5 (T5) & Functional analysis monographs, advanced QM, essential PhD papers\\
Cat.~A--F     & Full alphabetical catalogue: Popular [A], UG textbooks [B],
                Grad textbooks [C], Monographs [D], Primary papers [E], Reviews [F]\\
Quick-index   & Tier--category cross-index for rapid navigation\\
\bottomrule
\end{tabular}
\end{center}

%-------------------------------------------------------------
\subsection{How to Extend the Bibliography to a New Module or Discipline}
%-------------------------------------------------------------

\instructornote{%
The following protocol allows any instructor to build a compliant tiered bibliography
for a new module (e.g.\ a new Series~II or Series~III module not yet covered).
}

\begin{tcolorbox}[goldbox]
\textbf{Five-step bibliography construction protocol:}

\medskip
\begin{enumerate}[leftmargin=2em,itemsep=6pt]
  \item \textbf{Identify one [N] entry per Tier~1.}
        Choose a popular/narrative book that gives physical intuition for the
        module's core concept without formalism. If none exists, a well-written
        Wikipedia article or a Nobel lecture can substitute.

  \item \textbf{Identify two or three [T] entries for Tiers~2--3.}
        One standard course textbook (Tier~2) and one more rigorous text or
        bridge text (Tier~3). Annotate each with the specific chapters covering
        this module's lectures.

  \item \textbf{Identify the 2--3 primary papers [P] that every Tier~3--5 student
        must encounter.}
        These are typically: the paper that first introduced the concept, the paper
        that proved the central theorem, and (optionally) the most-cited modern
        review. Annotate with which lecture code they support and which WRQ/ORQ
        prompts they feed into.

  \item \textbf{Identify one or two [M] monographs for Tiers~4--5.}
        One standard graduate reference (Tier~4) and one advanced monograph
        or review survey (Tier~5). Annotate with chapters, not whole books.

  \item \textbf{Assign tier badges and star ratings.}
        Use the three-star system. Essential (\essential) means: a student at this
        tier \emph{must} read this to pass the Complex Problem Set (CPSU/CPSG)
        and Research Questions (WRQ/ORQ) for this module.
\end{enumerate}
\end{tcolorbox}

%-------------------------------------------------------------
\subsection{Relationship Between Bibliography and Assessment Artifacts}
%-------------------------------------------------------------

The tiered bibliography is not decorative --- it feeds directly into the assessment
artifacts defined in Section~\ref{sec:artifacts}. The table below makes the
connections explicit:

\begin{center}
\begin{tabularx}{\textwidth}{llX}
\toprule
\textbf{Artifact} & \textbf{Track} & \textbf{Bibliography connection}\\
\midrule
SCQU / CCQU & UG    & Students check answers against [N] and Tier~2 [T] entries\\
SPSU        & UG    & Problems drawn from exercises in starred Tier~2 [T] texts\\
CPSU        & UG    & Proofs derived from Tier~3 [T] and first [P] papers;
                      students cite the source of each theorem they use\\
SPJU        & UG    & Implementations reproduce a result from a starred [T] text\\
CPJU        & UG    & Report situates the student's comparison within the
                      Tier~2--3 textbook literature\\
SCQG / CCQG & Grad  & Precise statements verified against Tier~4--5 [T] and [M] entries\\
CPSG        & Grad  & Full proofs reference Tier~5 [M] sources (Reed--Simon, Rudin, etc.)\\
SPJG        & Grad  & Each of the 5 projects implements or verifies a result from
                      a specific [P] or [M] entry\\
CPJG        & Grad  & Each of the 5 complex projects requires a mini literature review
                      using Cat.~C--D entries\\
WRQ         & Res.  & Each of the 5 questions is directly seeded by a [P] entry;
                      the memo must cite at least 5 items from the bibliography\\
ORQ         & Res.  & Each of the 5 questions pushes beyond the bibliography;
                      students must identify at least 2 items \emph{not} in the
                      current catalogue and justify their inclusion\\
\bottomrule
\end{tabularx}
\end{center}

\medskip
\noindent The ORQ's requirement to find items \emph{not} in the bibliography is
deliberate: it trains the research skill of \emph{extending} a curated list, not
just consuming it.

%-------------------------------------------------------------
\subsection{Cross-Discipline Transfer: The Nonlinear Dynamics Example}
%-------------------------------------------------------------

The companion bibliography includes cross-reference entries from the
\textbf{Nonlinear Dynamics \& Chaos} (logistic map) course at each tier.
These serve two purposes:

\begin{itemize}[leftmargin=1.8em,itemsep=3pt]
  \item \textbf{Structural illustration:} They show concretely how the same
        five-tier architecture transfers to a completely different field.
        Gleick's \textit{Chaos} [T1] plays the same role as Feynman's \textit{QED} [T1];
        May (1976) [P, T3] plays the same role as Born (1926) [P, T3];
        de Melo \& van Strien [M, T5] plays the same role as Reed \& Simon [M, T5].

  \item \textbf{Interdisciplinary resonance:} Several quantum mechanics modules
        (especially Series~II Modules~II.7 and~II.5) touch phenomena --- ergodicity,
        invariant measures, spectral statistics --- that have direct nonlinear
        dynamics analogues. The cross-references make these connections visible
        to students at the right tier.
\end{itemize}

%-------------------------------------------------------------
\subsection{Maintenance and Version Control}
%-------------------------------------------------------------

\begin{tcolorbox}[goldbox]
\textbf{Recommended update cycle for the bibliography companion:}

\medskip
\begin{itemize}[leftmargin=1.5em,itemsep=4pt]
  \item \textbf{Annual review:} Check all URLs, update edition numbers, and scan
        arXiv/journal alerts for new papers that should replace or supplement
        existing [P] or [M] entries in Tiers~4--5.

  \item \textbf{Module release:} Whenever a new module is added to the programme,
        apply the five-step protocol (Section~3.6 above) to generate its
        bibliography entries and add them to the relevant tier sections.

  \item \textbf{Student feedback loop:} ORQ submissions frequently surface
        high-quality papers not in the catalogue. A nominated student each year
        can maintain a ``candidates'' list; the instructor reviews it and incorporates
        the best additions in the next annual update.

  \item \textbf{Version labelling:} The bibliography document carries a version
        number (e.g.\ \texttt{bib-v1.0}) separate from this PSM document. The
        two are linked but versioned independently so the bibliography can be
        updated without triggering a full PSM revision.
\end{itemize}
\end{tcolorbox}

%=============================================================
\end{document}
